\Needspace{18\baselineskip}
\section{Shallow plugging, including self-loop holes}
Fix a host $C:I\to O$, an edge $h$ to remove, and a replacement $H$ whose boundary
lists are exactly the input and output signature lists of $C$'s label at $h$.
All constructions in this section are in the namespace $\mathsf{Net.Replacement}$
unless explicitly stated otherwise.

\Needspace{24\baselineskip}
\subsection{Carriers and boundary identification}
\begin{definition}[Replacement carriers]
The new internal vertices are $C_V+H_V$, with disjoint tags providing freshness.
The new edges are $(C_E\setminus\{h\})+H_E$. The outer boundary remains $I,O$.
The surviving-edge subtype carries a proof of inequality to $h$, so the removed
edge cannot accidentally survive as an old edge.
\end{definition}
\lean{Graph/Plugging}{19}{19}{Net.Replacement.Vertex}
\lean{Graph/Plugging}{20}{20}{Net.Replacement.Edge}
The host embedding $j$ tags each old internal vertex on the left and fixes the
outer boundary. The insertion map $k$ tags a replacement internal vertex on the
right, sends replacement input $i$ to $j(s_h(i))$, and sends replacement output
$i$ to $j(t_h(i))$. Thus gluing is expressed by a function rather than a quotient
of all vertex names.
\lean{Graph/Plugging}{22}{22}{Net.Replacement.host}
\lean{Graph/Plugging}{24}{27}{Net.Replacement.insert}
The host embedding is injective by injectivity of the left sum tag and the port
map lemma. Its image is disjoint from freshly tagged replacement internals:
case analysis on a host vertex gives either a boundary constructor or an internal
left tag, neither equal to an internal right tag.
\lean{Graph/Plugging}{29}{30}{Net.Replacement.host_injective}
\lean{Graph/Plugging}{32}{33}{Net.Replacement.host_ne_internal}

\begin{lemma}[The two restricted insertion injections]
The insertion map is injective on vertices that are not outputs. It is also
injective on vertices that are not inputs.
\end{lemma}
\begin{proof}
For the first claim, exclude output cases using the hypotheses. Two internal
images agree only when their right-tagged names agree. An internal image cannot
agree with an input image, by freshness. Two input images agree only when their
host source ports agree, and source-port injectivity identifies their indices.
For the second claim, exclude input cases instead, use target-port injectivity
on two outputs, and use the same freshness argument on mixed cases.
\end{proof}
\lean{Graph/Plugging}{35}{50}{Net.Replacement.insert_source_injective}
\lean{Graph/Plugging}{52}{67}{Net.Replacement.insert_target_injective}
If a non-output replacement vertex has insertion image $j(v)$, it must be some
input $i$, with $s_h(i)=v$. Outputs are excluded by hypothesis, internals by
freshness, and in the input case host injectivity gives the equation. Dually, a
non-input replacement vertex mapping to $j(v)$ must be output $i$ with
$t_h(i)=v$. These inversion lemmas control mixed old/new incidence cases.
\lean{Graph/Plugging}{69}{75}{Net.Replacement.insert_source_host}
\lean{Graph/Plugging}{77}{83}{Net.Replacement.insert_target_host}

\Needspace{24\baselineskip}
\subsection{Preservation of incidence invariants}
Sources and targets of surviving host edges are their old lists mapped by $j$.
Those of inserted edges are the replacement lists mapped by $k$.
\lean{Graph/Plugging}{85}{87}{Net.Replacement.source}
\lean{Graph/Plugging}{89}{91}{Net.Replacement.target}
\begin{lemma}[Distinct incidence entries]
Every new source list and every new target list has no duplicates.
\end{lemma}
\begin{proof}
A surviving host list uses the injective map $j$. A replacement source list
contains no outputs, so the first restricted injection of $k$ applies to its
members. A replacement target list contains no inputs, so the second injection
applies. In all cases use preservation of distinctness under restricted injectivity.
\end{proof}
\lean{Graph/Plugging}{93}{98}{Net.Replacement.source_nodup}
\lean{Graph/Plugging}{100}{105}{Net.Replacement.target_nodup}

\begin{lemma}[Unique source and target incidence]
A vertex belongs to the sources of at most one new edge, and to the targets
of at most one new edge.
\end{lemma}
\begin{proof}
For source incidence, split on the two edge tags. Two host edges reduce, by
injectivity of $j$, to old source uniqueness. Two inserted edges reduce, by
restricted source injectivity of $k$, to source uniqueness in $H$. In the mixed
case the shared vertex is an inserted input, by the source inversion lemma;
its corresponding host vertex is therefore a source port of $h$. Old source
uniqueness identifies the surviving host edge with $h$, contradicting its subtype
certificate. The other mixed ordering is the same contradiction with the equality
reversed. For target incidence, repeat the four cases with target uniqueness,
restricted target injectivity, and the target inversion lemma: the mixed vertex
is now an inserted output and a target port of $h$.
\end{proof}
\lean{Graph/Plugging}{107}{136}{Net.Replacement.source_unique}
\lean{Graph/Plugging}{138}{167}{Net.Replacement.target_unique}
A vertex in a new source list is not an outer output. In the host case use the
old source property and preservation of boundary constructors. In the inserted
case an internal vertex remains internal, an input maps to a source port of $h$,
and an output is impossible in a replacement source list. For a new target list,
exchange source/input with target/output: inserted inputs are impossible and
inserted outputs map to host target ports, which cannot be outer inputs.
\lean{Graph/Plugging}{169}{180}{Net.Replacement.source_not_output}
\lean{Graph/Plugging}{182}{193}{Net.Replacement.target_not_input}

\begin{lemma}[Coverage of surviving host vertices]
For every host vertex $v$ that is not an output, $j(v)$ belongs to a new source
list. If $v$ is not an input, $j(v)$ belongs to a new target list.
\end{lemma}
\begin{proof}
Choose an old source edge by coverage in $C$. If it differs from $h$, retain it.
If it equals $h$, express $v$ as source port $i$; replacement input $i$ belongs
to a source edge by coverage in $H$, and its image is $j(v)$. For target coverage,
choose an old target edge. Retain it unless it is $h$; in that case express $v$
as target port $i$ and use coverage at replacement output $i$.
\end{proof}
\lean{Graph/Plugging}{195}{204}{Net.Replacement.source_host_covered}
\lean{Graph/Plugging}{206}{215}{Net.Replacement.target_host_covered}
Full source coverage now follows by splitting a new vertex into a host internal,
a fresh internal, an outer input, or an outer output. Host internals and inputs
use the preceding lemma. Fresh internals use replacement coverage and map the
witness edge into the right edge tag. Outer outputs contradict the required
non-output hypothesis. The reverse implication is the source exclusion already
proved. Full target coverage uses host internals and outer outputs, replacement
coverage for fresh internals, and excludes outer inputs; its reverse implication
is target exclusion.
\lean{Graph/Plugging}{217}{231}{Net.Replacement.source_coverage}
\lean{Graph/Plugging}{233}{247}{Net.Replacement.target_coverage}

\Needspace{24\baselineskip}
\subsection{Labels, typing, and the certified result}
New internal labels use the host label on left tags and replacement label on
right tags. Surviving edges keep their labels; inserted edges receive their
replacement labels.
\lean{Graph/Plugging}{249}{249}{Net.Replacement.internalLabel}
\lean{Graph/Plugging}{251}{253}{Net.Replacement.edgeLabel}
Labelling commutes with $j$, by the three port cases. It also commutes with $k$:
internal labels agree by definition, while input and output labels agree by the
host indexed-port typing lemmas and the matching replacement interface.
\lean{Graph/Plugging}{255}{257}{Net.Replacement.host_label}
\lean{Graph/Plugging}{259}{264}{Net.Replacement.insert_label}
For source typing, split the edge tag and compose the incidence map with the
vertex-label map. The commuting equation just proved reduces the result to
source typing in $C$ or $H$, respectively. For target typing the same composition
argument reduces to their target typing equations.
\lean{Graph/Plugging}{266}{276}{Net.Replacement.source_typed}
\lean{Graph/Plugging}{278}{288}{Net.Replacement.target_typed}

\begin{theorem}[Well-formed shallow plugging]
There is a net $C[h:=H]:I\to O$ with the carriers, labels, and incidence lists
above. It preserves the entire ordered external interface and all graph invariants,
even when $h$ has a vertex in both its source and target lists.
\end{theorem}
\begin{proof}
Enumerate the new internal carrier by the sum enumeration. Enumerate surviving
host edges by subset filtering, and concatenate the tagged replacement-edge
enumeration. Supply each of the distinctness, uniqueness, coverage, and typing
proofs established above as the corresponding net field. This assembles the
certified graph. Classical filtering of the surviving-edge subtype accounts for
the noncomputable qualifier.
\end{proof}
\lean{Graph/Plugging}{294}{312}{Net.plug}
There is deliberately no claim that $k$ is globally injective in this theorem.
If a host source port equals a host target port, its replacement input and output
are identified. The separate source and target injection arguments are sufficient
for every net invariant.
